\paragraph{Proofs - Global Markov Property }\hfill

The proof of the global Markov property follows similar arguments as used in \cite{Lau90, Ver93, Ric03, FM17, FM18, RERS17}, 
namely chaining the separoid axioms together in an inductive way.
The main difference here is that we never rely on the Symmetry property but instead use 
the left and right versions of the separoid axioms separately.


\begin{Thm}[Global Markov property for causal Bayesian networks]
\label{GMP-CCBN}
Consider a causal Bayesian network (with input and latent variables) with observational CADMG $G=(J,V,E,L)$ 
and observational Markov kernel $P(X_V|\doit(X_J))$.
Then for all $A, B, C \ins J \cup V$ (not-necessarily disjoint) we have the implication:
\[ A \Perp^d_G B \given C \qquad \implies \qquad X_A \Indep_{P(X_V|\doit(X_J))} X_B \given X_C.   \]
If one wants to make the implicit dependence on $J$ more explicit one can equivalently also write:
\[ A \Perp^d_G J \cup B \given C \qquad \implies \qquad X_A \Indep_{P(X_V|\doit(X_J))} X_J,X_{B} \given X_C.   \]
\begin{proof}
Because d-separation is preserved under marginalization:
\[ A \Perp^d_G B \given C \quad \iff \quad A \Perp^d_{G^+} B \given C,   \]
we can directly assume that we work with the causal Bayesian network without latent variables that marginalizes to the given one.
So w.l.o.g.\ $L=\emptyset$ and $G$ is a CDAG.\\
We then do induction by $\#V$. \\

%For this let $<$ be a topological order of $G$ where all elements of $J$ are ordered first.
%With this we have:
%\[ P(X_V|\doit(X_J)) = \bigotimes_{v \in V} P_v(X_v|\doit(X_{\Pa^G(v)})),\]
%where the product goes in reverse order of $<$. We will write $\Perp$ for $\Perp^d_G$ and $\Indep$ for 
%$\Indep_{P(X_V|\doit(X_J))}$ in the following.\\

0.) Induction start:  $V= \emptyset$. This means that $A,B,C \ins J$.
The assumption: 
\[A \Perp^d_G  B \given C,\] 
implies that we must have that $A \ins C$. Otherwise a trivial walk from $A \ins J$ to $J \cup B$ would be $C$-open.
Since $A,B,C \ins J$ we have the factorization:
\[ P(X_A,X_B,X_C|\doit(X_J)) = \underbrace{\bigotimes_{w \in A} \delta(X_w|X_w)}_{=:Q(X_A|X_C)} \otimes 
 \underbrace{\bigotimes_{w \in B} \delta(X_w|X_w) \otimes \bigotimes_{w \in C} \delta(X_w|X_w)}_{=P(X_B,X_C|\doit(X_J))}.  \]
 Because $A \ins C$ the Markov kernel $Q(X_A|X_C):= \bigotimes_{w \in A} \delta(X_w|X_w) $ 
 really is a Markov kernel from $\Xcal_C \dshto \Xcal_A$.
This already shows:
\[ X_A \Indep_{P(X_V|\doit(X_J))} X_B |X_C.\]

(IND): Induction assumption: 
The global Markov property holds for all causal Bayesian networks 
(with input variables, but without latent variables and without bi-directed edges) 
with $\#V < n$ (and arbitrary $J$).\\

1.) Now assume: $\#V =n>0$ and $A \Perp^d_G B \given C$.

Since $G$ is acyclic we can find a topological order $<$ for $G$ where the elements of $J$ are ordered first.  
Let $v \in V$ be its last element, which is thus  childless.  \\
Note that, since $\Ch^G(v) = \emptyset$, the marginalization $G^{\sm \{v\}}$ has no bi-directed edges and thus induces again 
a causal Bayesian network without latent variables with $\#V^{\sm \{v\}}=n-1 <n$.\\
Furthermore, we have the factorization:
\begin{align*} P(X_V|\doit(X_J)) &= P_v(X_v|\doit(X_{\Pa^G(v)})) \otimes 
    \underbrace{\bigotimes_{w \in \Pred^G_<(v) \sm J} P_w(X_w|\doit(X_{\Pa^G(w)}))}_{P(X_{\Pred_<^G(v) \sm J}|\doit(X_J))}.
%   &= P_v(X_v|\doit(X_{\Pa^G(v)})) \otimes  P(X_{\Pred_<^G(v)}|\doit(X_J)).
\end{align*}
This factorization implies that we already have the conditional independence:
\begin{align*}
    X_v \Indep_{P(X_V|\doit(X_J)) } X_{\Pred^G_<(v)} \given X_D,     \label{eq:pred}
\end{align*}
where we put $D:=\Pa^{G}(v)$.\\

In the following we will distinguish between 4 cases: 
\begin{enumerate}[label = \Alph*.)]
    \item $v \in A \sm C$,
    \item $v \in B \sm C$,
    \item $v \in C$,
    \item $ v \notin A \cup J \cup B \cup C$,
\end{enumerate}
Note that $v \in V$, thus $v \notin J$, which shows that the above cover all possible cases.\\
Further note that:
\[A \Perp^d_G  B \given C,\] 
implies that:
\[ A \cap ( J \cup B) \ins C.  \]
Otherwise a trivial walk from $A$ to $J \cup B$ would be $C$-open. 
This shows that $A\sm C$, $(J \cup B) \sm C$ and $C$ are pairwise disjoint.\\


Case D.):  $v \notin A \cup J \cup B \cup C$. Then we can marginalize out $v$ and use the equivalence: 
\[A \Perp^d_{G}  B \given C\quad \iff \quad A \Perp^d_{G^{\sm v}}  B \given C.\]
With $\#V^{\sm \{v\}}<n$ and induction (IND) we then get:
\[X_A \Indep_{P(X_V|\doit(X_J))} X_B \given X_C.\]
This shows the claim in case D.\\

Case A.):  $v \in A\sm C$. Then we can write:
\begin{align*}
    A &= A' \dcup (A \cap C) \dcup \{v\}, \\
    B &= B' \dcup (B \cap C),
\end{align*}
with some disjoint $A' \ins A \sm C$ and $B' \ins B \sm C$.
We then have the implications:
\begin{align*}
    A \Perp^d_G  B \given C & \xRightarrow{\text{Right Decomposition}}  & A \Perp^d_G B' \given C \\
    & \xRightarrow{\text{Left Decomposition}}  & A' \Perp^d_G  B' \given C \\
    & \xRightarrow{\text{marginalization, } v \notin A'\cup J \cup B' \cup C}  & A' \Perp^d_{G^{\sm \{v\}}}  B' \given C \\
 & \xRightarrow{\text{induction (IND)}}  & X_{A'} \Indep_{P(X_V|\doit(X_J))}  X_{B'} \given X_C. \label{eqn:bli} \tag{\#1}\\
\end{align*}
On the other hand we have with $D = \Pa^G(v)$:
\begin{align*}
 A \Perp^d_G  B \given C & \xRightarrow{\text{Right Decomposition, }B' \ins B}  & A \Perp^d_G   B' \given C   \\
  & \xRightarrow{\text{Left Weak Union, } A= A' \dcup (A \cap C) \dcup \{v\}}  & \{v\} \Perp^d_G  B' \given A' \dcup C    \\
 & \xRightarrow{(*), \text{ see below}}  & D  \Perp^d_G   B' \given A' \dcup C \\ 
 & \xRightarrow{\text{marginalization, } v \notin D \cup J \cup B' \cup A' \cup C}  & D \Perp^d_{G^{\sm \{v\}}}  B' \given A' \dcup C \\
& \xRightarrow{\text{induction (IND)}}  &   X_D  \Indep_{P(X_V|\doit(X_J))}  X_{B'} \given X_{A' \dcup C}\\
& \xRightarrow{A'\dcup C}  &   X_D  \Indep_{P(X_V|\doit(X_J))}  X_{B'} \given X_{A'}, X_C. \label{eqn:bla} \tag{\#2}
\end{align*}
(*) holds since every $(A' \dcup C)$-open walk $ w \sus \cdots$ from a  $w \in D=\Pa^G(v)$ to $J \cup B'$ 
extends to an $(A' \dcup C)$-open walk from $v$ to $J \cup B'$ via  $v \hut w \sus \cdots$, as $w$ stays a 
non-collider in the extended walk (not in $A' \dcup C$) and $v \notin A' \dcup C$.\\

As discussed above we also already have the conditional independence:
\[ X_v \Indep_{P(X_V|\doit(X_J))}  X_{\Pred^G_<(v)} \given X_D. \]
With this and $A' \dcup B' \dcup C \ins \Pred^G_<(v)$ we get the implications:
\begin{align*}
&&X_v \Indep_{P(X_V|\doit(X_J))}  X_{\Pred^G_<(v)} \given X_D\\ 
&\xRightarrow{\text{Right Decomposition}} & 
X_v \Indep_{P(X_V|\doit(X_J))}  X_{A'}, X_{B'}, X_C \given X_D \\
& \xRightarrow{\text{Right Weak Union}}  & 
X_v \Indep_{P(X_V|\doit(X_J))} X_{B'} \given X_{A'}, X_C, X_D \\ 
& \xRightarrow{\text{Left Contraction, }\eqref{eqn:bla}}  & X_v, X_D  \Indep_{P(X_V|\doit(X_J))}  X_{B'} \given X_{A'}, X_C \\
& \xRightarrow{\text{Left Decomposition}}  & X_v \Indep_{P(X_V|\doit(X_J))}  X_{B'} \given X_{A'}, X_C\\
& \xRightarrow{\text{Left Contraction, } \eqref{eqn:bli}}  &  X_{A'},X_v \Indep_{P(X_V|\doit(X_J))}  X_{B'} \given X_C\\
& \xRightarrow{X_J\text{-Inverted Right Decomposition}}  & X_{A'}, X_v \Indep_{P(X_V|\doit(X_J))}  X_J,X_{B'},X_C \given X_C\\
& \xRightarrow{\text{Right Decompositon, } B \ins B' \dcup C}  & X_{A'}, X_v \Indep_{P(X_V|\doit(X_J))}  X_B \given X_C.
 \label{eqn:bloop} \tag{\#3}
\end{align*}
By (Extended) Left Redundancy we have:
\[ X_{A'},X_v, X_C \Indep_{P(X_V|\doit(X_J))}  X_B \given X_{A'},X_v,X_C. \]
With this we get the implications:
\begin{align*}
&&X_{A'},X_v, X_C \Indep_{P(X_V|\doit(X_J))}  X_B \given X_{A'},X_v,X_C\\
& \xRightarrow{\text{Left Contraction, } \eqref{eqn:bloop}}  &  
X_{A'},X_v,X_{A'},X_v, X_C \Indep_{P(X_V|\doit(X_J))}  X_B \given X_C\\
& \xRightarrow{\text{Left Decomposition, } A \ins A' \dcup \{v\} \dcup C}  &    
X_A \Indep_{P(X_V|\doit(X_J))}  X_B \given X_C. 
\end{align*}
This shows the claim in case A.\\

%%%%%%%%
%
Case B.):  $v \in B\sm C$. Then we can write:
\begin{align*}
    A &= A' \dcup (A \cap C), \\
    B &= B' \dcup (B \cap C) \dcup \{v\},
\end{align*}
with some disjoint $A' \ins A \sm C$ and $B' \ins B \sm C$.

%Then $B= \{v\} \dot \cup B'$ and $A \Indep_{G^+}^\sigma \{v\},B',J \given C$ holds.
We then have the implications:
\begin{align*}
A  \Perp^d_G B \given C 
& \xRightarrow{\text{Left Decomposition}}  & A' \Perp^d_G  B \given C \\
& \xRightarrow{\text{Right Decomposition}}  & A' \Perp^d_G  B' \given C \\
 & \xRightarrow{\text{marginalization, } v \notin A'\cup J \cup B' \cup C}  & A' \Perp^d_{G^{\sm \{v\}}}  B' \given C \\
 & \xRightarrow{\text{induction (IND)}}  & X_{A'} \Indep_{P(X_V|\doit(X_J))}  X_{B'} \given X_C. \label{eqn:bli'} \tag{\#1'}
\end{align*}
Again with $D = \Pa^G(v)$ we get:
\begin{align*}
A  \Perp^d_G  B \given C  & \xRightarrow{\text{Left Decomposition}}  & A' \Perp^d_G  B \given C \\
 & \xRightarrow{\text{Right Decomposition}}  & A' \Perp^d_G  B' \cup \{v\} \given C \\
& \xRightarrow{\text{Right Weak Union}}  & A' \Perp^d_G  \{v\} \given B' \dcup C  \\
 & \xRightarrow{(\bullet), \text{ see below}}  & A'  \Perp^d_G  D \given B' \dcup C \\  
  & \xRightarrow{\text{marginalization, } v \notin A' \cup J \cup D \cup B' \cup C}  & A'  \Perp^d_{G^{\sm \{v\}}}  D \given B' \dcup C \\  
& \xRightarrow{\text{induction (IND)}}  &   X_{A'}  \Indep_{P(X_V|\doit(X_J))}  X_D \given X_{B' \dcup C}\\
& \xRightarrow{ B' \dcup C}  &   X_{A'}  \Indep_{P(X_V|\doit(X_J))}  X_D \given X_{B'}, X_C. \label{eqn:bla'} \tag{\#2'} 
\end{align*}
$(\bullet)$ holds since every $(B' \dcup C)$-open walk $\cdots \sus w$ from $A'$ to a $w \in J \cup D$ extends to 
a $(B' \dcup C)$-open walk from $A'$ to $J \cup \{v\}$, either because $w \in J$ or via $\cdots \sus w \tuh v$ if 
$w \in D=\Pa^G(v)$.
Note again that $w$ stays a non-collider in the extended walk (outside of $B'\dcup C$) and $v \notin B'\dcup C$.\\

As before we will use the following conditional independence:
\[ X_v \Indep_{P(X_V|\doit(X_J))}  X_{\Pred^G_<(v)} \given X_D. \]
With this and $A' \cup J \cup B' \cup C \ins \Pred^G_<(v)$ we get the implications:
\begin{align*}
   && X_v \Indep_{P(X_V|\doit(X_J))} X_{\Pred^G_<(v)} \given X_D \\
    &\xRightarrow{\text{Right Decomposition}} & X_v \Indep_{P(X_V|\doit(X_J))} X_{A'},X_{B'},X_C \given X_D\\
& \xRightarrow{\text{Right Weak Union}}  &  X_v \Indep_{P(X_V|\doit(X_J))} X_{A'} \given X_{B'},X_C,X_D\\
& \xRightarrow{\text{Flipped Left Cross Contraction, }\eqref{eqn:bla'}}  
& X_{A'}  \Indep_{P(X_V|\doit(X_J))} X_D,X_v \given X_{B'},X_C \\
& \xRightarrow{\text{Right Decomposition}}  
& X_{A'}  \Indep_{P(X_V|\doit(X_J))} X_v \given X_{B'},X_C \\
& \xRightarrow{\text{Right Contraction, }\eqref{eqn:bli'}}  
& X_{A'}  \Indep_{P(X_V|\doit(X_J))} X_{B'},X_v \given X_C \\
& \xRightarrow{\text{$X_J$-Inverted Right Decomposition}}  & X_{A'}  \Indep_{P(X_V|\doit(X_J))} X_J,X_{B'},X_v,X_C \given X_C \\
& \xRightarrow{\text{Right Decomposition, } B \ins B'\dcup \{v\} \dcup C}  & X_{A'}  \Indep_{P(X_V|\doit(X_J))} X_B \given X_C.  \label{eqn:bloop'} \tag{\#3'}
\end{align*}
By Redundancy we have:
\[ X_{A'}, X_C \Indep_{P(X_V|\doit(X_J))}  X_B \given X_{A'},X_C. \]
With this we get the implications:
\begin{align*}
&&X_{A'}, X_C \Indep_{P(X_V|\doit(X_J))}  X_B \given X_{A'},X_C\\
& \xRightarrow{\text{Left Contraction, } \eqref{eqn:bloop'}}  &  
X_{A'},X_{A'}, X_C \Indep_{P(X_V|\doit(X_J))}  X_B \given X_C.\\
& \xRightarrow{\text{Left Decomposition, } A \ins A' \dcup C}  &    
X_A \Indep_{P(X_V|\doit(X_J))}  X_B \given X_C. 
\end{align*}
This shows the claim in case B.\\


Case C.): $v \in C$.
 Then we can write:
\begin{align*}
    A &= A' \dcup (A \cap C), \\
    %J \cup B &= B' \dcup ((J \cup B) \cap C),\\
    B &= B' \dcup (B \cap C),\\
    C &= C' \dcup \{v\},
\end{align*}
with some pairwise disjoint $A' \ins A \sm C$, 
%$B' \ins (J \cup B) \sm C$, 
$B' \ins B \sm C$ and
$C' \ins C$.

We then get the implications.
\begin{align*}
A  \Perp^d_G  B \given C 
& \xRightarrow{\text{Left Decomposition}}  & A' \Perp^d_G  B \given C \\
& \xRightarrow{\text{Right Decomposition}}  & A' \Perp^d_G  B' \given C \\
& \xRightarrow{C = C' \dcup \{v\}}  & A' \Perp^d_G  B' \given C'\dcup\{v\} \\
\end{align*}


%Now assume the case when $v \in C$, so $C= \{v \} \dot \cup C'$.
%Using the same arguments as in A.0 before we can w.l.o.g.\ assume that $A$ is disjoint from $J$ (this is to avoid making assumptions about standardness of $\Xcal_j$, $j \in J$). %\Patrick{$v  \notin J$ !!}
%Then we have: $A \Indep_{G^+}^\sigma B,J \given \{v\},C'$.

We now claim that:
\[A' \Perp^d_G  B' \given C'\dcup\{v\}\]
implies that one of the following statements holds:
\[  A' \dcup \{v\} \Perp^d_G  B' \given C' \qquad \lor \qquad A' \Perp^d_G B' \dcup \{v\} \given C'. \]
Assume the contrary:
\[  A' \dcup \{v\} \nPerp^d_G  B' \given C' \qquad \land \qquad A' \nPerp^d_G B' \dcup \{v\} \given C'. \]
So there exist shortest $C'$-open walks $\pi_1$ and $\pi_2$ in $G$ such that all colliders are in $C'$:
\[ \pi_1:\quad  A' \cup \{v\}\ni u_0 \sus \cdots \sus u_k \in J \cup B',  \]
and:
\[ \pi_2:\quad  A' \ni w_0 \sus \cdots \sus w_m \in J \cup (B' \dcup \{v\}).  \]
So all non-colliders of $\pi_1$ and $\pi_2$ are outside of $C'$.
Since we consider shortest walks and $v \notin C'$ at most an end node of $\pi_1$ and $\pi_2$ could be equal to $v$.
Otherwise one could shorten the walk.\\
Then note that $v \notin A'$ and $v \notin J \cup B'$, thus: $u_k \neq v$ and $w_0 \neq v$.\\
If now $\pi_i$ does not contain $v$ as an (end) node, then $\pi_i$ would be $(C'\dcup \{v\})$-open,
which is a contradiction to the assumption:
\[A' \Perp^d_G  B' \given C'\dcup\{v\}.\]
So we can assume that the other end nodes equal $v$, i.e.: $u_0=v$ and $w_m =v$.\\
Furthermore, both $\pi_1$ and $\pi_2$ are non-trivial walks, since $u_0 \neq u_k$ and $w_0 \neq w_m$.
%Otherwise $v=u_k \in J \cup B'$ or $v=w_0 \in A'$, both of which are not true.
%So $k,m \ge 1$.\\
Since $v$ is childless and $k,m \ge 1$ we have that the $\pi_i$ are of the forms:
\[ \pi_1:\quad   v \hut u_1 \sus \cdots \sus u_k,  \]
and:
\[ \pi_2:\quad  w_0 \sus \cdots \sus w_{m-1} \tuh v,  \]
with $u_1, w_{m-1} \in D=\Pa^G(v)$. Then the following walk:
\[ A' \ni  w_0 \sus \cdots \sus w_{m-1} \tuh v \hut u_1 \sus \cdots \sus u_k \in J \cup B',  \]
is a $(C' \dcup \{v\})$-open walk from $A'$ to $J \cup B'$, in contradiction to:
\[A' \Perp^d_G B' \given C'\dcup\{v\}.\]
So the claim:
\[  A' \dcup \{v\} \Perp^d_G  B' \given C' \qquad \lor \qquad A' \Perp^d_G B' \dcup \{v\} \given C', \]
must be true. 
So we reduced case C to case A or case B, which then imply:
\[ X_A,X_v \Indep_{P(X_V|\doit(X_J))}  X_B \given X_{C'} \qquad \lor \qquad X_A \Indep_{P(X_V|\doit(X_J))}  X_B,X_v \given X_{C'}.\]
If we apply Left Weak Union to the left and Right Weak Union to the right we get:
\[ X_A \Indep_{P(X_V|\doit(X_J))}  X_B \given X_{C'},X_v,\]
which implies:
\[ X_A \Indep_{P(X_V|\doit(X_J))}  X_B \given X_C.\]
This shows the claim in case C.
\end{proof}

\end{Thm}

